\documentclass[12pt]{article}

\usepackage{amsmath}
\usepackage{amssymb}
\usepackage[margin=1in]{geometry}
\usepackage{hyperref}
\usepackage{xcolor}
\usepackage{booktabs}
\usepackage{array}
\usepackage{enumitem}
\usepackage{fancyhdr}

% Set up headers and footers
\pagestyle{fancy}
\fancyhf{}
\rhead{GED Math Answer Key}
\lhead{Number Sense and Operations}
\cfoot{\thepage}

% Document metadata
\title{\Huge GED Math Answer Key \\ \Large Number Sense and Operations (30 Questions)}
\author{Comprehensive Solutions and Explanations}
\date{June 2026}

% Define colors
\definecolor{correctgreen}{RGB}{0, 128, 0}
\definecolor{lightgreen}{RGB}{200, 255, 200}
\definecolor{lightblue}{RGB}{200, 230, 255}

\begin{document}

% Title Page
\maketitle
\thispagestyle{empty}
\vspace{2cm}

\begin{center}
\Large \textbf{Complete Answer Key with Detailed Solutions} \\[0.5cm]
\textit{Three Difficulty Levels: Basic, Intermediate, and Advanced} \\[2cm]
\normalsize
This document provides comprehensive solutions for all 30 GED Math Practice questions, featuring:\\
$\bullet$ Step-by-step problem-solving strategies \\
$\bullet$ Mathematical concept explanations \\
$\bullet$ Common mistakes to avoid \\
$\bullet$ Alternative solution methods \\
$\bullet$ Real-world applications \\[2cm]
Designed to support student learning and understanding of fundamental mathematics skills.
\end{center}

\newpage

% Answer Summary Table
\section*{Answer Summary}

\begin{table}[h]
\centering
\small
\begin{tabular}{|c|c|c|c|c|c|c|}
\hline
\textbf{Q} & \textbf{Answer} & \textbf{Q} & \textbf{Answer} & \textbf{Q} & \textbf{Answer} \\
\hline
1 & \colorbox{lightgreen}{A} & 11 & \colorbox{lightgreen}{B} & 21 & \colorbox{lightgreen}{B} \\
\hline
2 & \colorbox{lightgreen}{B} & 12 & \colorbox{lightgreen}{C} & 22 & \colorbox{lightgreen}{B} \\
\hline
3 & \colorbox{lightgreen}{B} & 13 & \colorbox{lightgreen}{C} & 23 & \colorbox{lightgreen}{C} \\
\hline
4 & \colorbox{lightgreen}{A} & 14 & \colorbox{lightgreen}{C} & 24 & \colorbox{lightgreen}{D} \\
\hline
5 & \colorbox{lightgreen}{A} & 15 & \colorbox{lightgreen}{D} & 25 & \colorbox{lightgreen}{D} \\
\hline
6 & \colorbox{lightgreen}{A} & 16 & \colorbox{lightgreen}{A} & 26 & \colorbox{lightgreen}{C} \\
\hline
7 & \colorbox{lightgreen}{A} & 17 & \colorbox{lightgreen}{A} & 27 & \colorbox{lightgreen}{D} \\
\hline
8 & \colorbox{lightgreen}{C} & 18 & \colorbox{lightgreen}{C} & 28 & \colorbox{lightgreen}{C} \\
\hline
9 & \colorbox{lightgreen}{C} & 19 & \colorbox{lightgreen}{A} & 29 & \colorbox{lightgreen}{B} \\
\hline
10 & \colorbox{lightgreen}{C} & 20 & \colorbox{lightgreen}{A} & 30 & \colorbox{lightgreen}{B} \\
\hline
\end{tabular}
\end{table}

\newpage

% Table of Contents
\tableofcontents
\newpage

% ==================== BASIC LEVEL ====================
\section{Basic Level: Fundamental Operations (Questions 1--10)}

\subsection{Question 1: Order of Operations}

\textbf{Correct Answer: } \colorbox{lightgreen}{A}

\subsubsection*{Problem}
Simplify: $8 + 3 \times 5 - 2$

\subsubsection*{Solution}

To solve this problem, we must follow the \textbf{Order of Operations}, often remembered by the acronym \textit{PEMDAS} (Parentheses, Exponents, Multiplication/Division, Addition/Subtraction).

\begin{enumerate}
\item \textbf{Step 1:} Identify the operations: addition, multiplication, and subtraction
\item \textbf{Step 2:} Perform multiplication first: $3 \times 5 = 15$
\item \textbf{Step 3:} Rewrite the expression: $8 + 15 - 2$
\item \textbf{Step 4:} Perform addition and subtraction from left to right:
  \begin{align*}
  8 + 15 &= 23 \\
  23 - 2 &= 21
  \end{align*}
\end{enumerate}

\textbf{Answer: } $21$

\subsubsection*{Mathematical Concept}
The Order of Operations is a standard convention that ensures everyone solves expressions the same way. Multiplication and division are performed before addition and subtraction (unless parentheses override this).

\subsubsection*{Common Mistakes to Avoid}
\begin{itemize}
\item Working strictly from left to right: $8 + 3 = 11$, then $11 \times 5 = 55$, etc. (incorrect)
\item Forgetting to multiply before adding: treating it as $(8 + 3) \times 5 - 2$ (incorrect)
\end{itemize}

\subsubsection*{Alternative Verification}
Using standard notation: $(8) + (3 \times 5) - (2) = 8 + 15 - 2 = 21$ 

\subsubsection*{Real-World Application}
Order of operations is crucial in budgeting and finance. For example: Starting with \$8, earning \$3 per hour for 5 hours, then spending \$2 gives you: $8 + 3(5) - 2 = \$21$.

---

\subsection{Question 2: Adding Fractions}

\textbf{Correct Answer: } \colorbox{lightgreen}{B}

\subsubsection*{Problem}
Add: $\dfrac{1}{4} + \dfrac{1}{3}$

\subsubsection*{Solution}

To add fractions, we need a common denominator.

\begin{enumerate}
\item \textbf{Step 1:} Find the least common denominator (LCD) of 4 and 3
  \begin{itemize}
  \item Multiples of 4: 4, 8, 12, 16, ...
  \item Multiples of 3: 3, 6, 9, 12, ...
  \item LCD = 12
  \end{itemize}
\item \textbf{Step 2:} Convert each fraction to have denominator 12:
  \begin{align*}
  \frac{1}{4} &= \frac{1 \times 3}{4 \times 3} = \frac{3}{12} \\
  \frac{1}{3} &= \frac{1 \times 4}{3 \times 4} = \frac{4}{12}
  \end{align*}
\item \textbf{Step 3:} Add the fractions:
  \[
  \frac{3}{12} + \frac{4}{12} = \frac{3 + 4}{12} = \frac{7}{12}
  \]
\end{enumerate}

\textbf{Answer: } $\dfrac{7}{12}$

\subsubsection*{Mathematical Concept}
When adding fractions, we can only add the numerators if the denominators are the same. Finding a common denominator allows us to rewrite fractions in equivalent forms.

\subsubsection*{Common Mistakes to Avoid}
\begin{itemize}
\item Adding numerators and denominators separately: $\frac{1+1}{4+3} = \frac{2}{7}$ (incorrect)
\item Using an incorrect common denominator: Using 7 instead of 12 (incorrect)
\end{itemize}

\subsubsection*{Alternative Verification}
Convert to decimals: $\frac{1}{4} = 0.25$ and $\frac{1}{3} \approx 0.333$, so $0.25 + 0.333 \approx 0.583$. Check: $\frac{7}{12} \approx 0.583$ 

\subsubsection*{Real-World Application}
Cooking requires adding fractions: If a recipe calls for $\frac{1}{4}$ cup flour and $\frac{1}{3}$ cup sugar, you need a total of $\frac{7}{12}$ cup of dry ingredients.

---

\subsection{Question 3: Subtracting Decimals}

\textbf{Correct Answer: } \colorbox{lightgreen}{B}

\subsubsection*{Problem}
Subtract: $8.92 - 5.45$

\subsubsection*{Solution}

To subtract decimals, align the decimal points and subtract as with whole numbers.

\begin{enumerate}
\item \textbf{Step 1:} Align the decimal points:
  \[
  \begin{array}{r}
  8.92 \\
  -5.45 \\
  \hline
  \end{array}
  \]
\item \textbf{Step 2:} Subtract from right to left:
  \begin{itemize}
  \item Hundredths place: $2 - 5$ requires borrowing. $12 - 5 = 7$
  \item Tenths place: $8 - 1 - 4 = 3$ (after borrowing 1)
  \item Ones place: $8 - 5 = 3$
  \end{itemize}
\item \textbf{Step 3:} Write the result:
  \[
  8.92 - 5.45 = 3.47
  \]
\end{enumerate}

\textbf{Answer: } $3.47$

\subsubsection*{Mathematical Concept}
Decimal subtraction follows the same regrouping rules as whole number subtraction. The decimal point remains in the same position.

\subsubsection*{Common Mistakes to Avoid}
\begin{itemize}
\item Misaligning decimal points: $8.92 - 5.45$ calculated as $892 - 545$ (incorrect)
\item Forgetting to borrow when necessary, leading to negative intermediate results
\end{itemize}

\subsubsection*{Alternative Verification}
Addition check: $3.47 + 5.45 = 8.92$ 

\subsubsection*{Real-World Application}
Shopping and making change: If an item costs \$5.45 and you pay with \$8.92, your change is $8.92 - 5.45 = \$3.47$.

---

\subsection{Question 4: Converting Fractions to Decimals}

\textbf{Correct Answer: } \colorbox{lightgreen}{A}

\subsubsection*{Problem}
Convert to a decimal: $\dfrac{3}{8}$

\subsubsection*{Solution}

To convert a fraction to a decimal, divide the numerator by the denominator.

\begin{enumerate}
\item \textbf{Step 1:} Set up the division: $3 \div 8$
\item \textbf{Step 2:} Since $3 < 8$, the result is less than 1. Begin with $0.$
\item \textbf{Step 3:} Perform long division:
  \begin{align*}
  3 \div 8 &= 0 \text{ remainder } 3 \\
  30 \div 8 &= 3 \text{ remainder } 6 \quad \text{(first decimal place)} \\
  60 \div 8 &= 7 \text{ remainder } 4 \quad \text{(second decimal place)} \\
  40 \div 8 &= 5 \text{ remainder } 0 \quad \text{(third decimal place)}
  \end{align*}
\item \textbf{Step 4:} Combine the results: $0.375$
\end{enumerate}

\textbf{Answer: } $0.375$

\subsubsection*{Mathematical Concept}
A fraction can be expressed as a division problem. The fraction bar represents division, so $\frac{3}{8} = 3 \div 8$.

\subsubsection*{Common Mistakes to Avoid}
\begin{itemize}
\item Computing $8 \div 3$ instead of $3 \div 8$ (incorrect)
\item Stopping the division prematurely, giving an approximation instead of the exact value
\end{itemize}

\subsubsection*{Alternative Verification}
Multiply back: $0.375 \times 8 = 3$ 

\subsubsection*{Real-World Application}
Engineering and measurements: $\frac{3}{8}$ inch on a ruler is exactly $0.375$ inches, important for precision work.

---

\subsection{Question 5: Finding Percentages}

\textbf{Correct Answer: } \colorbox{lightgreen}{A}

\subsubsection*{Problem}
15 is what percent of 60?

\subsubsection*{Solution}

Use the percentage formula: $\text{Percentage} = \dfrac{\text{Part}}{\text{Whole}} \times 100\%$

\begin{enumerate}
\item \textbf{Step 1:} Identify the part and whole:
  \begin{itemize}
  \item Part = 15
  \item Whole = 60
  \end{itemize}
\item \textbf{Step 2:} Apply the formula:
  \[
  \text{Percentage} = \frac{15}{60} \times 100\%
  \]
\item \textbf{Step 3:} Simplify the fraction:
  \[
  \frac{15}{60} = \frac{1}{4} = 0.25
  \]
\item \textbf{Step 4:} Multiply by 100\%:
  \[
  0.25 \times 100\% = 25\%
  \]
\end{enumerate}

\textbf{Answer: } $25\%$

\subsubsection*{Mathematical Concept}
A percentage expresses a number as a fraction of 100. The formula relates a part to a whole, giving us the percentage.

\subsubsection*{Common Mistakes to Avoid}
\begin{itemize}
\item Reversing part and whole: $\frac{60}{15} \times 100\% = 400\%$ (incorrect)
\item Forgetting to multiply by 100 to convert to a percentage
\end{itemize}

\subsubsection*{Alternative Verification}
Check: $25\% \text{ of } 60 = 0.25 \times 60 = 15$ 

\subsubsection*{Real-World Application}
Test scores: If you answer 15 questions correctly out of 60 total, you scored 25%.

---

\subsection{Question 6: Multiplying Fractions}

\textbf{Correct Answer: } \colorbox{lightgreen}{A}

\subsubsection*{Problem}
Multiply: $\dfrac{2}{3} \times \dfrac{1}{5}$

\subsubsection*{Solution}

To multiply fractions, multiply numerators together and denominators together, then simplify.

\begin{enumerate}
\item \textbf{Step 1:} Multiply numerators: $2 \times 1 = 2$
\item \textbf{Step 2:} Multiply denominators: $3 \times 5 = 15$
\item \textbf{Step 3:} Write the result:
  \[
  \frac{2}{3} \times \frac{1}{5} = \frac{2 \times 1}{3 \times 5} = \frac{2}{15}
  \]
\item \textbf{Step 4:} Check if simplification is possible:
  \begin{itemize}
  \item Factors of 2: 1, 2
  \item Factors of 15: 1, 3, 5, 15
  \item GCD = 1, so the fraction is already in simplest form
  \end{itemize}
\end{enumerate}

\textbf{Answer: } $\dfrac{2}{15}$

\subsubsection*{Mathematical Concept}
Unlike addition and subtraction, multiplying fractions does not require a common denominator. Multiply straight across and simplify the result.

\subsubsection*{Common Mistakes to Avoid}
\begin{itemize}
\item Finding a common denominator (not needed for multiplication)
\item Forgetting to simplify the final answer
\item Multiplying denominators incorrectly
\end{itemize}

\subsubsection*{Alternative Verification}
Convert to decimals: $\frac{2}{3} \approx 0.667$ and $\frac{1}{5} = 0.2$, so $0.667 \times 0.2 \approx 0.133$. Check: $\frac{2}{15} \approx 0.133$ 

\subsubsection*{Real-World Application}
Cooking: If you need $\frac{2}{3}$ of a $\frac{1}{5}$ cup ingredient, you need $\frac{2}{15}$ cup.

---

\subsection{Question 7: Dividing Fractions}

\textbf{Correct Answer: } \colorbox{lightgreen}{A}

\subsubsection*{Problem}
Divide: $\dfrac{2}{3} \div \dfrac{3}{4}$

\subsubsection*{Solution}

To divide by a fraction, multiply by its reciprocal (the fraction flipped).

\begin{enumerate}
\item \textbf{Step 1:} Identify the divisor: $\frac{3}{4}$
\item \textbf{Step 2:} Find its reciprocal: $\frac{4}{3}$
\item \textbf{Step 3:} Change division to multiplication:
  \[
  \frac{2}{3} \div \frac{3}{4} = \frac{2}{3} \times \frac{4}{3}
  \]
\item \textbf{Step 4:} Multiply:
  \[
  \frac{2}{3} \times \frac{4}{3} = \frac{2 \times 4}{3 \times 3} = \frac{8}{9}
  \]
\item \textbf{Step 5:} Simplify (if needed):
  \begin{itemize}
  \item GCD(8, 9) = 1, so $\frac{8}{9}$ is in simplest form
  \end{itemize}
\end{enumerate}

\textbf{Answer: } $\dfrac{8}{9}$

\subsubsection*{Mathematical Concept}
Dividing by a number is equivalent to multiplying by its reciprocal. For fractions, the reciprocal of $\frac{a}{b}$ is $\frac{b}{a}$.

\subsubsection*{Common Mistakes to Avoid}
\begin{itemize}
\item Flipping the wrong fraction: Using $\frac{3}{2}$ instead of $\frac{4}{3}$ (incorrect)
\item Performing division instead of multiplication after flipping
\item Forgetting the step-by-step process and trying to divide directly
\end{itemize}

\subsubsection*{Alternative Verification}
Check: $\frac{8}{9} \times \frac{3}{4} = \frac{24}{36} = \frac{2}{3}$ 

\subsubsection*{Real-World Application}
Splitting portions: If you have $\frac{2}{3}$ of a pizza and want to divide it into $\frac{3}{4}$-sized slices, you get $\frac{8}{9}$ slices.

---

\subsection{Question 8: Exponents}

\textbf{Correct Answer: } \colorbox{lightgreen}{C}

\subsubsection*{Problem}
Evaluate: $2^5$

\subsubsection*{Solution}

An exponent indicates how many times to multiply a base by itself.

\begin{enumerate}
\item \textbf{Step 1:} Identify the base and exponent:
  \begin{itemize}
  \item Base = 2
  \item Exponent = 5
  \end{itemize}
\item \textbf{Step 2:} Expand the expression:
  \[
  2^5 = 2 \times 2 \times 2 \times 2 \times 2
  \]
\item \textbf{Step 3:} Multiply step by step:
  \begin{align*}
  2 \times 2 &= 4 \\
  4 \times 2 &= 8 \\
  8 \times 2 &= 16 \\
  16 \times 2 &= 32
  \end{align*}
\end{enumerate}

\textbf{Answer: } $32$

\subsubsection*{Mathematical Concept}
Exponents are a shorthand for repeated multiplication. The number of times the base is multiplied is determined by the exponent.

\subsubsection*{Common Mistakes to Avoid}
\begin{itemize}
\item Multiplying base times exponent: $2 \times 5 = 10$ (incorrect)
\item Miscounting the number of multiplications
\item Confusing $2^5$ with $5^2$
\end{itemize}

\subsubsection*{Alternative Verification}
Powers of 2 sequence: $2^1=2$, $2^2=4$, $2^3=8$, $2^4=16$, $2^5=32$ 

\subsubsection*{Real-World Application}
Computing: Binary systems use powers of 2. $2^5 = 32$ represents the number of possible combinations with 5 binary digits.

---

\subsection{Question 9: Square Roots}

\textbf{Correct Answer: } \colorbox{lightgreen}{C}

\subsubsection*{Problem}
Find: $\sqrt{64}$

\subsubsection*{Solution}

The square root of a number is the value that, when multiplied by itself, equals that number.

\begin{enumerate}
\item \textbf{Step 1:} Ask: What number times itself equals 64?
\item \textbf{Step 2:} Test values:
  \begin{itemize}
  \item $7 \times 7 = 49$ (too small)
  \item $8 \times 8 = 64$ (correct!)
  \item $9 \times 9 = 81$ (too large)
  \end{itemize}
\item \textbf{Step 3:} Verify: $8 \times 8 = 64$ 
\end{enumerate}

\textbf{Answer: } $8$

\subsubsection*{Mathematical Concept}
A square root is the inverse of squaring. If $n^2 = 64$, then $\sqrt{64} = n$. The symbol $\sqrt{\phantom{x}}$ represents the principal (positive) square root.

\subsubsection*{Common Mistakes to Avoid}
\begin{itemize}
\item Finding half the number: $\sqrt{64} \neq 32$ (incorrect)
\item Confusing with division by 2
\item Forgetting that both $+8$ and $-8$ square to 64, but $\sqrt{64}$ refers to the positive root
\end{itemize}

\subsubsection*{Alternative Verification}
Check: $8^2 = 8 \times 8 = 64$ 

\subsubsection*{Real-World Application}
Geometry: If a square room has an area of 64 square feet, each side is $\sqrt{64} = 8$ feet.

---

\subsection{Question 10: Absolute Value}

\textbf{Correct Answer: } \colorbox{lightgreen}{C}

\subsubsection*{Problem}
Evaluate: $|-5|$

\subsubsection*{Solution}

The absolute value of a number is its distance from zero on the number line, always positive.

\begin{enumerate}
\item \textbf{Step 1:} Identify the number: $-5$
\item \textbf{Step 2:} Find its distance from 0:
  \begin{itemize}
  \item The number $-5$ is 5 units to the left of 0
  \item Distance is always positive
  \end{itemize}
\item \textbf{Step 3:} Apply the absolute value:
  \[
  |-5| = 5
  \]
\end{enumerate}

\textbf{Answer: } $5$

\subsubsection*{Mathematical Concept}
Absolute value $|x|$ gives the magnitude of a number without regard to its sign. For any real number $n$:
\[
|n| = \begin{cases} n & \text{if } n \geq 0 \\ -n & \text{if } n < 0 \end{cases}
\]

\subsubsection*{Common Mistakes to Avoid}
\begin{itemize}
\item Keeping the negative sign: $|-5| \neq -5$ (incorrect)
\item Confusing absolute value with negation
\item Not recognizing that $|5| = |-5| = 5$
\end{itemize}

\subsubsection*{Alternative Verification}
Both $5$ and $-5$ are distance 5 from zero: $|5| = 5$ and $|-5| = 5$ 

\subsubsection*{Real-World Application}
Temperature change: A change of $-5°\text{F}$ (decrease) has an absolute value of $5°\text{F}$ in magnitude.

\newpage

% ==================== INTERMEDIATE LEVEL ====================
\section{Intermediate Level: Multi-Step Operations (Questions 11--20)}

\subsection{Question 11: Complex Order of Operations}

\textbf{Correct Answer: } \colorbox{lightgreen}{B}

\subsubsection*{Problem}
Simplify: $24 \div (4 + 2) - 3 + 1$

\subsubsection*{Solution}

Follow PEMDAS, paying special attention to the parentheses.

\begin{enumerate}
\item \textbf{Step 1:} Parentheses first: $4 + 2 = 6$
\item \textbf{Step 2:} Rewrite: $24 \div 6 - 3 + 1$
\item \textbf{Step 3:} Division: $24 \div 6 = 4$
\item \textbf{Step 4:} Rewrite: $4 - 3 + 1$
\item \textbf{Step 5:} Subtraction and addition (left to right):
  \begin{align*}
  4 - 3 &= 1 \\
  1 + 1 &= 2
  \end{align*}
\end{enumerate}

\textbf{Answer: } $2$

\subsubsection*{Mathematical Concept}
Parentheses must be evaluated first in the order of operations. After parentheses are resolved, continue with the standard order: exponents, multiplication/division, addition/subtraction.

\subsubsection*{Common Mistakes to Avoid}
\begin{itemize}
\item Performing operations without parentheses first: $24 \div 4 = 6$, then $6 + 2 = 8$, etc. (incorrect)
\item Working strictly left to right without considering order of operations
\item Dividing 24 by 4 only, ignoring the 2 inside the parentheses
\end{itemize}

\subsubsection*{Alternative Verification}
Double-check each step: $(6 \text{ groups of } 4) \div 6 = 4$; $4 - 3 + 1 = 2$ 

\subsubsection*{Real-World Application}
Recipe adjustments: If 24 tablespoons must be divided equally among $(4 + 2)$ servings, minus 3 tablespoons reserved, plus 1 extra, you use 2 tablespoons per serving.

---

\subsection{Question 12: Converting Percentages (Discount)}

\textbf{Correct Answer: } \colorbox{lightgreen}{C}

\subsubsection*{Problem}
A \$50 shirt is on sale for 25\% off. What is the final price?

\subsubsection*{Solution}

Calculate the discount amount, then subtract from the original price.

\begin{enumerate}
\item \textbf{Step 1:} Calculate the discount:
  \[
  \text{Discount} = 25\% \times \$50 = 0.25 \times \$50 = \$12.50
  \]
\item \textbf{Step 2:} Subtract from original price:
  \[
  \text{Final Price} = \$50 - \$12.50 = \$37.50
  \]
\end{enumerate}

\textbf{Answer: } \$37.50

\subsubsection*{Mathematical Concept}
A percentage discount reduces the original price. The final price is the original price minus the discount amount, or equivalently, the original price times $(1 - \text{discount percentage})$.

\subsubsection*{Common Mistakes to Avoid}
\begin{itemize}
\item Finding 25\% of 50 correctly (\$12.50) but then adding instead of subtracting
\item Calculating 25\% as 0.025 instead of 0.25
\item Forgetting to subtract the discount amount
\end{itemize}

\subsubsection*{Alternative Verification}
Method 2: If 25\% is discounted, you pay $100\% - 25\% = 75\%$ of the original price:
\[
75\% \times \$50 = 0.75 \times \$50 = \$37.50 \quad \checkmark
\]

\subsubsection*{Real-World Application}
Shopping: A 25\% discount on a \$50 shirt saves you \$12.50, making your purchase \$37.50.

---

\subsection{Question 13: Ratio and Proportion}

\textbf{Correct Answer: } \colorbox{lightgreen}{C}

\subsubsection*{Problem}
A recipe calls for flour and sugar in a ratio of 3:2. If you use 8 cups of sugar, how much flour is needed?

\subsubsection*{Solution}

Set up a proportion using the given ratio.

\begin{enumerate}
\item \textbf{Step 1:} Write the ratio of flour to sugar: $\frac{\text{flour}}{\text{sugar}} = \frac{3}{2}$
\item \textbf{Step 2:} Set up a proportion with the unknown:
  \[
  \frac{3}{2} = \frac{x}{8}
  \]
  where $x$ is the amount of flour needed.
\item \textbf{Step 3:} Cross-multiply:
  \[
  3 \times 8 = 2 \times x
  \]
\item \textbf{Step 4:} Solve for $x$:
  \begin{align*}
  24 &= 2x \\
  x &= \frac{24}{2} = 12
  \end{align*}
\item \textbf{Step 5:} Check: $\frac{12}{8} = \frac{3}{2}$ 
\end{enumerate}

\textbf{Answer: } 12 cups of flour

\subsubsection*{Mathematical Concept}
A ratio compares two quantities. A proportion states that two ratios are equal. When one part of a proportion changes, we can find the corresponding change in the other part.

\subsubsection*{Common Mistakes to Avoid}
\begin{itemize}
\item Setting up the ratio backwards: $\frac{2}{3}$ instead of $\frac{3}{2}$ (incorrect)
\item Solving $2x = 8$ instead of $2x = 24$
\item Forgetting to cross-multiply correctly
\end{itemize}

\subsubsection*{Alternative Verification}
Scaling factor: $\frac{8}{2} = 4$. So flour is $3 \times 4 = 12$ cups 

\subsubsection*{Real-World Application}
Cooking: Recipes use ratios to maintain proportions. A 3:2 ratio of flour to sugar scales consistently whether you double or halve the recipe.

---

\subsection{Question 14: Multi-Step Percentage}

\textbf{Correct Answer: } \colorbox{lightgreen}{C}

\subsubsection*{Problem}
A house costs \$200,000. It increases in value by 20\%, then decreases by 10\%. What is the final value?

\subsubsection*{Solution}

Apply percentage changes sequentially.

\begin{enumerate}
\item \textbf{Step 1:} Calculate the 20\% increase:
  \[
  \text{Increase} = 20\% \times \$200,000 = 0.20 \times \$200,000 = \$40,000
  \]
\item \textbf{Step 2:} New value after increase:
  \[
  \$200,000 + \$40,000 = \$240,000
  \]
\item \textbf{Step 3:} Calculate the 10\% decrease on the new value:
  \[
  \text{Decrease} = 10\% \times \$240,000 = 0.10 \times \$240,000 = \$24,000
  \]
\item \textbf{Step 4:} Final value after decrease:
  \[
  \$240,000 - \$24,000 = \$216,000
  \]
\end{enumerate}

\textbf{Answer: } \$216,000

\subsubsection*{Mathematical Concept}
When applying consecutive percentage changes, each change is applied to the result of the previous step, not the original amount. Order matters in multiplication, but sequential percentage applications are sequential operations.

\subsubsection*{Common Mistakes to Avoid}
\begin{itemize}
\item Adding 20\% and subtracting 10\% from the original: $200,000 + 40,000 - 20,000 = 220,000$ (incorrect; you'd be subtracting 10\% of the original, not the new value)
\item Calculating the 10\% decrease on the original \$200,000 instead of \$240,000
\item Assuming a 20\% increase and 10\% decrease nets to a 10\% overall increase
\end{itemize}

\subsubsection*{Alternative Verification}
Multiplier method: $200,000 \times 1.20 \times 0.90 = 200,000 \times 1.08 = 216,000$ 

\subsubsection*{Real-World Application}
Real estate: Property values fluctuate. A house gaining 20\% then losing 10\% ends up ahead overall (8\% net increase).

---

\subsection{Question 15: Fraction to Percent}

\textbf{Correct Answer: } \colorbox{lightgreen}{D}

\subsubsection*{Problem}
Convert to a percentage: $\dfrac{3}{5}$

\subsubsection*{Solution}

Convert the fraction to a decimal, then multiply by 100\%.

\begin{enumerate}
\item \textbf{Step 1:} Divide to get a decimal: $3 \div 5 = 0.6$
\item \textbf{Step 2:} Multiply by 100\%:
  \[
  0.6 \times 100\% = 60\%
  \]
\end{enumerate}

\textbf{Answer: } $60\%$

\subsubsection*{Mathematical Concept}
To convert a fraction to a percentage, divide the numerator by the denominator to get a decimal, then multiply by 100\% (or move the decimal point two places to the right).

\subsubsection*{Common Mistakes to Avoid}
\begin{itemize}
\item Dividing the wrong way: $5 \div 3$ instead of $3 \div 5$ (incorrect)
\item Multiplying by 100 without understanding it's converting to a percentage scale
\item Writing the answer as 0.6\% instead of 60\%
\end{itemize}

\subsubsection*{Alternative Verification}
Check: $60\% = \frac{60}{100} = \frac{3}{5}$ 

\subsubsection*{Real-World Application}
Grades: If you answer 3 out of 5 questions correctly, you score 60\%.

---

\subsection{Question 16: Negative Numbers}

\textbf{Correct Answer: } \colorbox{lightgreen}{A}

\subsubsection*{Problem}
Subtract: $-3 - 5$

\subsubsection*{Solution}

Subtracting is equivalent to adding the opposite.

\begin{enumerate}
\item \textbf{Step 1:} Rewrite subtraction as adding the opposite:
  \[
  -3 - 5 = -3 + (-5)
  \]
\item \textbf{Step 2:} Add two negative numbers (move left on the number line):
  \[
  -3 + (-5) = -(3 + 5) = -8
  \]
\end{enumerate}

\textbf{Answer: } $-8$

\subsubsection*{Mathematical Concept}
Subtracting a positive number is the same as adding a negative number. When adding two negative numbers, add their absolute values and apply the negative sign to the result.

\subsubsection*{Common Mistakes to Avoid}
\begin{itemize}
\item Computing $-3 - 5 = 2$ (wrong sign)
\item Thinking $-3 - 5 = -3 + 5 = 2$ (incorrect operation reversal)
\item Losing track of which number is positive and which is negative
\end{itemize}

\subsubsection*{Alternative Verification}
Number line: Start at $-3$, move 5 units left (negative direction) to reach $-8$ 

\subsubsection*{Real-World Application}
Bank accounts: Starting at $-\$3$ (in debt) and losing another $\$5$ results in $-\$8$ debt.

---

\subsection{Question 17: Multiplying Negatives}

\textbf{Correct Answer: } \colorbox{lightgreen}{A}

\subsubsection*{Problem}
Multiply: $3 \times (-4)$

\subsubsection*{Solution}

A positive number times a negative number yields a negative result.

\begin{enumerate}
\item \textbf{Step 1:} Recognize the signs: positive $\times$ negative
\item \textbf{Step 2:} Multiply the absolute values: $3 \times 4 = 12$
\item \textbf{Step 3:} Apply the sign: negative
  \[
  3 \times (-4) = -12
  \]
\end{enumerate}

\textbf{Answer: } $-12$

\subsubsection*{Mathematical Concept}
Sign rules for multiplication:
\begin{itemize}
\item Positive $\times$ Positive = Positive
\item Negative $\times$ Negative = Positive
\item Positive $\times$ Negative = Negative
\item Negative $\times$ Positive = Negative
\end{itemize}

\subsubsection*{Common Mistakes to Avoid}
\begin{itemize}
\item Forgetting the negative sign and getting 12
\item Mistakenly thinking positive $\times$ negative = positive
\item Confusing multiplication with addition of negatives
\end{itemize}

\subsubsection*{Alternative Verification}
Interpretation: Three groups of $-4$ items each: $(-4) + (-4) + (-4) = -12$ 

\subsubsection*{Real-World Application}
Temperature: If the temperature drops 4 degrees per hour for 3 hours, the total change is $3 \times (-4) = -12$ degrees.

---

\subsection{Question 18: Multiple Fractions}

\textbf{Correct Answer: } \colorbox{lightgreen}{C}

\subsubsection*{Problem}
Add: $\dfrac{1}{2} + \dfrac{1}{6}$

\subsubsection*{Solution}

Find a common denominator and add.

\begin{enumerate}
\item \textbf{Step 1:} Find the LCD of 2 and 6:
  \begin{itemize}
  \item Multiples of 2: 2, 4, 6, 8, ...
  \item Multiples of 6: 6, 12, 18, ...
  \item LCD = 6
  \end{itemize}
\item \textbf{Step 2:} Convert fractions to have denominator 6:
  \begin{align*}
  \frac{1}{2} &= \frac{1 \times 3}{2 \times 3} = \frac{3}{6} \\
  \frac{1}{6} &= \frac{1}{6}
  \end{align*}
\item \textbf{Step 3:} Add:
  \[
  \frac{3}{6} + \frac{1}{6} = \frac{4}{6} = \frac{2}{3}
  \]
\end{enumerate}

\textbf{Answer: } $\dfrac{2}{3}$

\subsubsection*{Mathematical Concept}
To add fractions with different denominators, find the LCD, convert each fraction, then add the numerators. Always simplify the final answer.

\subsubsection*{Common Mistakes to Avoid}
\begin{itemize}
\item Using 12 as the LCD when 6 works (inefficient but not incorrect if properly simplified)
\item Forgetting to simplify: leaving the answer as $\frac{4}{6}$ instead of $\frac{2}{3}$
\item Adding denominators: $\frac{1}{2} + \frac{1}{6} = \frac{1}{8}$ (incorrect)
\end{itemize}

\subsubsection*{Alternative Verification}
Convert to decimals: $\frac{1}{2} = 0.5$ and $\frac{1}{6} \approx 0.167$, so $0.5 + 0.167 \approx 0.667$. Check: $\frac{2}{3} \approx 0.667$ 

\subsubsection*{Real-World Application}
Cooking: $\frac{1}{2}$ cup plus $\frac{1}{6}$ cup equals $\frac{2}{3}$ cup of an ingredient.

---

\subsection{Question 19: Tax Calculation}

\textbf{Correct Answer: } \colorbox{lightgreen}{A}

\subsubsection*{Problem}
A book costs \$58.50. With 12\% sales tax, what is the total cost?

\subsubsection*{Solution}

Calculate the tax amount and add it to the original price.

\begin{enumerate}
\item \textbf{Step 1:} Calculate 12\% tax:
  \[
  \text{Tax} = 12\% \times \$58.50 = 0.12 \times \$58.50 = \$7.02
  \]
\item \textbf{Step 2:} Add tax to original price:
  \[
  \text{Total} = \$58.50 + \$7.02 = \$65.52
  \]
\end{enumerate}

\textbf{Answer: } \$65.52

\subsubsection*{Mathematical Concept}
Sales tax is added to the base price. The total cost is the original price plus the tax amount, or equivalently, the original price times $(1 + \text{tax rate})$.

\subsubsection*{Common Mistakes to Avoid}
\begin{itemize}
\item Calculating 12\% of 58.50 incorrectly
\item Subtracting tax instead of adding
\item Forgetting to add the tax to the base price
\item Rounding errors in decimal arithmetic
\end{itemize}

\subsubsection*{Alternative Verification}
Method 2: $\$58.50 \times 1.12 = \$65.52$ 

\subsubsection*{Real-World Application}
Shopping: When you buy a \$58.50 book, 12\% sales tax adds \$7.02, making your total \$65.52.

---

\subsection{Question 20: Comparing Forms}

\textbf{Correct Answer: } \colorbox{lightgreen}{A}

\subsubsection*{Problem}
Show that $\dfrac{3}{4} = 0.75 = 75\%$

\subsubsection*{Solution}

Convert the fraction through each form to verify equivalence.

\begin{enumerate}
\item \textbf{Step 1:} Convert fraction to decimal:
  \[
  \frac{3}{4} = 3 \div 4 = 0.75
  \]
\item \textbf{Step 2:} Convert decimal to percentage:
  \[
  0.75 \times 100\% = 75\%
  \]
\item \textbf{Step 3:} Verify all three forms represent the same value:
  \[
  \frac{3}{4} = 0.75 = 75\%
  \]
\end{enumerate}

\textbf{Answer: } All three forms are equivalent

\subsubsection*{Mathematical Concept}
Fractions, decimals, and percentages are different representations of the same quantity. Being fluent in converting between forms is essential for mathematical flexibility.

\subsubsection*{Common Mistakes to Avoid}
\begin{itemize}
\item Miscalculating the decimal: $3 \div 4 \neq 0.34$ (incorrect)
\item Forgetting to multiply by 100 when converting decimal to percentage
\item Assuming fractions and percentages are always different representations
\end{itemize}

\subsubsection*{Alternative Verification}
Work backwards: $75\% = \frac{75}{100} = \frac{3}{4}$ (after simplifying by dividing by 25) 

\subsubsection*{Real-World Application}
Grades: A score of $\frac{3}{4}$ on a test is the same as 0.75 or 75\%, all expressing the same performance level.

\newpage

% ==================== ADVANCED LEVEL ====================
\section{Advanced Level: Complex Problem-Solving (Questions 21--30)}

\subsection{Question 21: Complex Fractions}

\textbf{Correct Answer: } \colorbox{lightgreen}{B}

\subsubsection*{Problem}
Add: $\dfrac{5}{8} + \dfrac{1}{6}$

\subsubsection*{Solution}

Find the LCD and add the fractions.

\begin{enumerate}
\item \textbf{Step 1:} Find the LCD of 8 and 6:
  \begin{itemize}
  \item Prime factorization of 8: $2^3$
  \item Prime factorization of 6: $2 \times 3$
  \item LCD = $2^3 \times 3 = 24$
  \end{itemize}
\item \textbf{Step 2:} Convert to common denominator:
  \begin{align*}
  \frac{5}{8} &= \frac{5 \times 3}{8 \times 3} = \frac{15}{24} \\
  \frac{1}{6} &= \frac{1 \times 4}{6 \times 4} = \frac{4}{24}
  \end{align*}
\item \textbf{Step 3:} Add:
  \[
  \frac{15}{24} + \frac{4}{24} = \frac{19}{24}
  \]
\item \textbf{Step 4:} Check if simplification is needed:
  \begin{itemize}
  \item Factors of 19: 1, 19 (prime)
  \item Factors of 24: 1, 2, 3, 4, 6, 8, 12, 24
  \item GCD = 1, so $\frac{19}{24}$ is in simplest form
  \end{itemize}
\end{enumerate}

\textbf{Answer: } $\dfrac{19}{24}$

\subsubsection*{Mathematical Concept}
When denominators are coprime (share no common factors other than 1), the LCD is their product. When they share factors, use prime factorization to find the LCD efficiently.

\subsubsection*{Common Mistakes to Avoid}
\begin{itemize}
\item Using 48 as the LCD instead of 24 (using the product instead of the true LCD)
\item Miscalculating conversion factors
\item Incorrectly identifying whether the result can be simplified
\end{itemize}

\subsubsection*{Alternative Verification}
Decimal check: $\frac{5}{8} = 0.625$ and $\frac{1}{6} \approx 0.167$, so sum $\approx 0.792$. Check: $\frac{19}{24} \approx 0.792$ 

\subsubsection*{Real-World Application}
Sewing: Combining $\frac{5}{8}$ yard of fabric with $\frac{1}{6}$ yard requires a total of $\frac{19}{24}$ yard.

---

\subsection{Question 22: Exponents}

\textbf{Correct Answer: } \colorbox{lightgreen}{B}

\subsubsection*{Problem}
Evaluate: $5^2 + 4^2$

\subsubsection*{Solution}

Calculate each square, then add.

\begin{enumerate}
\item \textbf{Step 1:} Calculate $5^2$:
  \[
  5^2 = 5 \times 5 = 25
  \]
\item \textbf{Step 2:} Calculate $4^2$:
  \[
  4^2 = 4 \times 4 = 16
  \]
\item \textbf{Step 3:} Add:
  \[
  25 + 16 = 41
  \]
\end{enumerate}

\textbf{Answer: } $41$

\subsubsection*{Mathematical Concept}
Exponents must be calculated before addition. $(5 + 4)^2 = 81$ is different from $5^2 + 4^2 = 41$. This demonstrates that exponents do not distribute over addition.

\subsubsection*{Common Mistakes to Avoid}
\begin{itemize}
\item Adding first: $(5 + 4)^2 = 9^2 = 81$ (incorrect)
\item Multiplying base and exponent: $5 \times 2 + 4 \times 2 = 18$ (incorrect)
\item Computing $(5^2 + 4)^2$ or other order-of-operations errors
\end{itemize}

\subsubsection*{Alternative Verification}
This is a Pythagorean triple component: $3^2 + 4^2 = 5^2$ (famous), but $5^2 + 4^2 = 25 + 16 = 41$ 

\subsubsection*{Real-World Application}
Construction: If you have a 5-unit and 4-unit board, their combined square footage is $5^2 + 4^2 = 41$ square units.

---

\subsection{Question 23: Square Roots}

\textbf{Correct Answer: } \colorbox{lightgreen}{C}

\subsubsection*{Problem}
A square room has an area of 169 square feet. What is the length of each side?

\subsubsection*{Solution}

Use the formula for the area of a square and solve for the side length.

\begin{enumerate}
\item \textbf{Step 1:} Set up the equation:
  \[
  \text{Area} = s^2 = 169
  \]
  where $s$ is the side length.
\item \textbf{Step 2:} Solve by taking the square root:
  \[
  s = \sqrt{169}
  \]
\item \textbf{Step 3:} Find the square root (recognizing perfect squares):
  \begin{itemize}
  \item $12^2 = 144$ (too small)
  \item $13^2 = 169$ (correct!)
  \end{itemize}
\item \textbf{Step 4:} State the answer:
  \[
  s = 13 \text{ feet}
  \]
\end{enumerate}

\textbf{Answer: } 13 feet

\subsubsection*{Mathematical Concept}
For a square with area $A$, the side length is $s = \sqrt{A}$. Square roots are the inverse operation of squaring.

\subsubsection*{Common Mistakes to Avoid}
\begin{itemize}
\item Dividing by 2: $169 \div 2 = 84.5$ (incorrect)
\item Forgetting units in the answer
\item Approximating when an exact answer exists
\end{itemize}

\subsubsection*{Alternative Verification}
Check: $13 \times 13 = 169$ 

\subsubsection*{Real-World Application}
Interior design: A square room with 169 square feet of floor space has sides of 13 feet each.

---

\subsection{Question 24: Compound Percent}

\textbf{Correct Answer: } \colorbox{lightgreen}{D}

\subsubsection*{Problem}
An investment of \$100 increases by 8\% in the first year and then by 7\% in the second year. What is the final value?

\subsubsection*{Solution}

Apply sequential percentage increases.

\begin{enumerate}
\item \textbf{Step 1:} Calculate value after first year (8\% increase):
  \[
  \text{Value}_1 = \$100 \times 1.08 = \$108
  \]
\item \textbf{Step 2:} Calculate value after second year (7\% increase on the new value):
  \[
  \text{Value}_2 = \$108 \times 1.07 = \$115.56
  \]
\end{enumerate}

\textbf{Answer: } \$115.56

\subsubsection*{Mathematical Concept}
Compound growth applies successive percentage increases to the accumulated amount. The formula is:
\[
\text{Final Value} = \text{Initial} \times (1 + r_1) \times (1 + r_2) \times \ldots
\]
where $r_i$ represents each rate of change.

\subsubsection*{Common Mistakes to Avoid}
\begin{itemize}
\item Adding percentages: \$100 + 8\% + 7\% = \$115 (ignoring the time-based compounding)
\item Calculating the 7\% increase on the original \$100 instead of on \$108
\item Misunderstanding that the second year's growth applies to the first year's result
\end{itemize}

\subsubsection*{Alternative Verification}
Direct calculation: $100 \times 1.08 \times 1.07 = 100 \times 1.1556 = \$115.56$ 

\subsubsection*{Real-World Application}
Investment growth: An initial investment of \$100 growing at 8\% annually and then 7\% reaches \$115.56 after two years.

---

\subsection{Question 25: Advanced Order of Operations}

\textbf{Correct Answer: } \colorbox{lightgreen}{D}

\subsubsection*{Problem}
Simplify: $\left(\dfrac{1}{2} + \dfrac{1}{4}\right) \times 8 - 2$

\subsubsection*{Solution}

Work through parentheses, then multiplication, then subtraction.

\begin{enumerate}
\item \textbf{Step 1:} Simplify inside parentheses. Find LCD of 2 and 4:
  \begin{align*}
  \frac{1}{2} &= \frac{2}{4} \\
  \frac{2}{4} + \frac{1}{4} &= \frac{3}{4}
  \end{align*}
\item \textbf{Step 2:} Rewrite: $\frac{3}{4} \times 8 - 2$
\item \textbf{Step 3:} Multiply:
  \[
  \frac{3}{4} \times 8 = \frac{3 \times 8}{4} = \frac{24}{4} = 6
  \]
\item \textbf{Step 4:} Rewrite: $6 - 2$
\item \textbf{Step 5:} Subtract:
  \[
  6 - 2 = 4
  \]
\end{enumerate}

\textbf{Answer: } $4$

\textbf{Note:} The provided answer key indicates the answer should be 10, but based on the problem statement and standard order of operations, the correct answer is 4. However, if the problem intended $\left(\frac{1}{2} + \frac{1}{4}\right) \times 8 + 2$, the answer would be $6 + 2 = 8$. If the problem is $\left(\frac{1}{2} + \frac{1}{4}\right) \times (8 - 2)$, the answer would be $\frac{3}{4} \times 6 = 4.5$. Given the provided answer is D = 10, there may be a transcription discrepancy. For pedagogical purposes, this solution demonstrates the correct application of order of operations.

\subsubsection*{Mathematical Concept}
Parentheses create groupings that must be evaluated first. Only after parentheses are resolved do we proceed with multiplication and division (left to right), then addition and subtraction (left to right).

\subsubsection*{Common Mistakes to Avoid}
\begin{itemize}
\item Multiplying before adding inside parentheses
\item Forgetting the parentheses entirely
\item Subtracting before multiplying
\end{itemize}

\subsubsection*{Alternative Verification}
Step-by-step check: $(\frac{1}{2} + \frac{1}{4}) = \frac{3}{4}$; $\frac{3}{4} \times 8 = 6$; $6 - 2 = 4$ 

\subsubsection*{Real-World Application}
Recipes with multiple operations: Combine ingredients (parentheses), multiply the resulting mixture by a scaling factor, then adjust the final amount.

---

\subsection{Question 26: Complex Percentage}

\textbf{Correct Answer: } \colorbox{lightgreen}{C}

\subsubsection*{Problem}
An \$80 item is marked up by 50\%, then discounted by 20\%. What is the final price?

\subsubsection*{Solution}

Apply the markup, then apply the discount to the marked-up price.

\begin{enumerate}
\item \textbf{Step 1:} Calculate 50\% markup:
  \[
  \text{Markup} = 50\% \times \$80 = 0.50 \times \$80 = \$40
  \]
\item \textbf{Step 2:} Price after markup:
  \[
  \$80 + \$40 = \$120
  \]
\item \textbf{Step 3:} Calculate 20\% discount on the marked-up price:
  \[
  \text{Discount} = 20\% \times \$120 = 0.20 \times \$120 = \$24
  \]
\item \textbf{Step 4:} Final price after discount:
  \[
  \$120 - \$24 = \$96
  \]
\end{enumerate}

\textbf{Answer: } \$96

\subsubsection*{Mathematical Concept}
Sequential operations must be applied in order. A markup increases the price first, then a discount applies to that new price. The combined effect is a net increase, since the 20\% discount applies to the larger marked-up value.

\subsubsection*{Common Mistakes to Avoid}
\begin{itemize}
\item Computing the discount on the original \$80: $20\% \times 80 = \$16$, giving $120 - 16 = \$104$ (incorrect)
\item Adding and subtracting percentages: $+50\% - 20\% = +30\%$ applied to \$80 = \$104 (incorrect; doesn't account for sequential application)
\item Reversing the order (discount then markup)
\end{itemize}

\subsubsection*{Alternative Verification}
Multiplier method: $\$80 \times 1.50 \times 0.80 = \$80 \times 1.20 = \$96$ 

\subsubsection*{Real-World Application}
Retail: A store marks up inventory 50\%, then runs a 20\% discount sale. Customers benefit, but the store still profits (12\% net markup).

---

\subsection{Question 27: Fraction/Decimal/Percent Equivalence}

\textbf{Correct Answer: } \colorbox{lightgreen}{D}

\subsubsection*{Problem}
Which is NOT equivalent to 0.2?
\begin{itemize}
\item A) $\dfrac{1}{5}$
\item B) 20\%
\item C) 2\%
\item D) $\dfrac{1}{4}$
\end{itemize}

\subsubsection*{Solution}

Convert each option to a common form (decimal) and compare.

\begin{enumerate}
\item \textbf{Step 1:} Convert option A: $\frac{1}{5} = 1 \div 5 = 0.2$  (equivalent)
\item \textbf{Step 2:} Convert option B: $20\% = \frac{20}{100} = 0.20 = 0.2$  (equivalent)
\item \textbf{Step 3:} Convert option C: $2\% = \frac{2}{100} = 0.02$  (NOT equivalent)
\item \textbf{Step 4:} Convert option D: $\frac{1}{4} = 1 \div 4 = 0.25$  (NOT equivalent)
\end{enumerate}

\subsubsection*{Conclusion}
Both C (2\%) and D ($\frac{1}{4}$) are not equivalent to 0.2. The answer key specifies D as the answer; both would be correct.

\textbf{Answer: } D (though C is also not equivalent)

\subsubsection*{Mathematical Concept}
Equivalent values must have the same decimal representation. $\frac{1}{5} = 0.2$ and $20\% = 0.2$, but $\frac{1}{4} = 0.25$ and $2\% = 0.02$, both different from 0.2.

\subsubsection*{Common Mistakes to Avoid}
\begin{itemize}
\item Confusing 2\% with 20\%
\item Miscalculating $\frac{1}{4} = 0.2$ when it actually equals 0.25
\item Not converting to a common form before comparing
\end{itemize}

\subsubsection*{Alternative Verification}
Work from 0.2: $0.2 = \frac{1}{5} = 20\%$. Compare: $\frac{1}{4} = 0.25 \neq 0.2$ and $2\% = 0.02 \neq 0.2$ 

\subsubsection*{Real-World Application}
Data analysis: Distinguishing between equivalent representations (20\% vs. 2\%) is critical for correctly interpreting statistics and surveys.

---

\subsection{Question 28: Negative Number Operations}

\textbf{Correct Answer: } \colorbox{lightgreen}{C}

\subsubsection*{Problem}
Evaluate: $(-3) \times (-2) + (-5)$

\subsubsection*{Solution}

Apply multiplication first (order of operations), then addition.

\begin{enumerate}
\item \textbf{Step 1:} Identify the operations: multiplication and addition
\item \textbf{Step 2:} Perform multiplication first:
  \[
  (-3) \times (-2) = 6 \quad \text{(negative times negative = positive)}
  \]
\item \textbf{Step 3:} Rewrite: $6 + (-5)$
\item \textbf{Step 4:} Add:
  \[
  6 + (-5) = 6 - 5 = 1
  \]
\end{enumerate}

\textbf{Answer: } $1$

\subsubsection*{Mathematical Concept}
The sign rules for multiplication are critical:
\begin{itemize}
\item Negative $\times$ Negative = Positive
\item Negative $\times$ Positive = Negative
\item Positive $\times$ Negative = Negative
\item Positive $\times$ Positive = Positive
\end{itemize}
After multiplication yields a positive result, adding a negative number is equivalent to subtraction.

\subsubsection*{Common Mistakes to Avoid}
\begin{itemize}
\item Computing $(-3) \times (-2) = -6$ (incorrect sign)
\item Forgetting to follow order of operations (multiplying before adding)
\item Adding incorrectly: $6 + (-5) \neq 11$ or other errors
\end{itemize}

\subsubsection*{Alternative Verification}
Rearrange: $(-3) \times (-2) + (-5) = 6 - 5 = 1$ 

\subsubsection*{Real-World Application}
Accounting: Negative income (debt) times negative time (past) yields positive value; then subtract current expenses.

---

\subsection{Question 29: Proportional Reasoning}

\textbf{Correct Answer: } \colorbox{lightgreen}{B}

\subsubsection*{Problem}
If 5 workers can complete a project in 12 days, how many days will it take 10 workers to complete the same project (assuming constant productivity)?

\subsubsection*{Solution}

Use inverse proportional reasoning: more workers means fewer days needed.

\begin{enumerate}
\item \textbf{Step 1:} Recognize the inverse relationship:
  \[
  \text{Workers} \times \text{Days} = \text{Constant (total work)}
  \]
\item \textbf{Step 2:} Calculate the constant:
  \[
  5 \times 12 = 60 \text{ worker-days}
  \]
\item \textbf{Step 3:} Set up proportion with 10 workers:
  \[
  10 \times d = 60
  \]
  where $d$ is the number of days.
\item \textbf{Step 4:} Solve for $d$:
  \[
  d = \frac{60}{10} = 6 \text{ days}
  \]
\end{enumerate}

\textbf{Answer: } 6 days

\subsubsection*{Mathematical Concept}
Inverse proportionality states that as one variable increases, the other decreases such that their product remains constant. Doubling the workers halves the time needed.

\subsubsection*{Common Mistakes to Avoid}
\begin{itemize}
\item Using direct proportion (doubling workers, doubling days): $10 \times 24 = 240$ (incorrect)
\item Forgetting to account for the inverse relationship
\item Arithmetic errors in calculating worker-days or final division
\end{itemize}

\subsubsection*{Alternative Verification}
Logic check: 5 workers take 12 days; 10 workers (double) should take half the time: $12 \div 2 = 6$ days 

\subsubsection*{Real-World Application}
Project management: More team members working on a project reduces the timeline proportionally (assuming independent, non-overlapping tasks).

---

\subsection{Question 30: Simple Interest}

\textbf{Correct Answer: } \colorbox{lightgreen}{B}

\subsubsection*{Problem}
You invest \$500 in a savings account with a 4\% annual interest rate. How much will you have after one year?

\subsubsection*{Solution}

Use the simple interest formula to find the interest earned, then add to principal.

\begin{enumerate}
\item \textbf{Step 1:} Identify the values:
  \begin{itemize}
  \item Principal $(P) = \$500$
  \item Rate $(r) = 4\% = 0.04$
  \item Time $(t) = 1 \text{ year}$
  \end{itemize}
\item \textbf{Step 2:} Apply the simple interest formula:
  \[
  I = P \times r \times t = \$500 \times 0.04 \times 1 = \$20
  \]
\item \textbf{Step 3:} Calculate total amount (principal + interest):
  \[
  A = P + I = \$500 + \$20 = \$520
  \]
\end{enumerate}

\textbf{Answer: } \$520

\subsubsection*{Mathematical Concept}
Simple interest is calculated as a percentage of the principal applied only once (or once per specified period). The formula is:
\[
I = P \times r \times t, \quad A = P + I = P(1 + rt)
\]
This contrasts with compound interest, where interest is calculated on accumulated interest.

\subsubsection*{Common Mistakes to Avoid}
\begin{itemize}
\item Calculating 4\% of 500 as 0.04 instead of 0.04 (correct but easy to misunderstand)
\item Forgetting to add interest back to principal
\item Confusing simple interest with compound interest
\item Miscalculating the percentage: $\$500 \times 0.4 = \$200$ (using 40\% instead of 4\%)
\end{itemize}

\subsubsection*{Alternative Verification}
Percentage method: $4\% \text{ of } \$500 = 0.04 \times 500 = \$20$; total = $\$500 + \$20 = \$520$ 

\subsubsection*{Real-World Application}
Banking: A \$500 deposit in a savings account at 4\% annual interest earns \$20 in interest over one year, giving you a total of \$520.

\newpage

\appendix

\section*{Summary and Key Concepts}

This answer key covers fundamental to advanced topics in arithmetic and financial mathematics:

\begin{enumerate}
\item \textbf{Order of Operations (PEMDAS):} Parentheses, Exponents, Multiplication/Division, Addition/Subtraction
\item \textbf{Fractions:} Adding, subtracting, multiplying, and dividing with common denominators
\item \textbf{Decimals:} Arithmetic operations and conversions to/from fractions and percentages
\item \textbf{Percentages:} Finding percentages, calculating discounts, markups, and tax
\item \textbf{Exponents and Roots:} Repeated multiplication and inverse operations
\item \textbf{Negative Numbers:} Operations with negative values and sign rules
\item \textbf{Ratios and Proportions:} Direct and inverse proportional relationships
\item \textbf{Financial Mathematics:} Percentages, taxes, discounts, simple interest, and investment growth
\end{enumerate}

Mastery of these concepts builds a strong foundation for higher-level mathematics and practical problem-solving.

\end{document}